\chapter{Kubernetes: Cluster Anatomy and the Reconciliation Loop}
\label{ch:cluster}

\section{The one idea everything else hangs on}

You do not tell Kubernetes to \emph{do} things; you tell it what should
\emph{be true}. You write \emph{desired state} (``a Deployment named
course-service with 1 replica of image X'') into its database via the
API server; \emph{controllers} run reconciliation loops forever
comparing desired with actual and acting on the difference. Delete a
pod: its controller notices actual $\neq$ desired and makes another.
Nobody replays your commands --- the state converges.

This explains behaviors that otherwise seem strange:

\begin{itemize}
  \item \code{kubectl apply} is a merge into desired state, not a
  script run --- apply the same file twice, nothing happens the second
  time (idempotence --- the property the pipeline leans on in
  Chapter~23).
  \item There is no ``startup order.'' Everything starts at once;
  services that need absent dependencies crash and are restarted until
  the world is ready. Crash-looping \emph{during} convergence is normal
  operation, not failure.
  \item Deleting from a manifest file does not delete from the cluster
  --- apply adds and updates, it does not prune. Removing the old
  course-service Ingress took an explicit \code{kubectl delete}.
\end{itemize}

\section{The moving parts}

\begin{figure}[h]
\centering
\begin{tikzpicture}[node distance=4mm and 10mm]
  \node[boxdark] (api) {API server};
  \node[boxfill, right=of api] (etcd) {state store\\\scriptsize (etcd / sqlite)};
  \node[box, below left=6mm and -2mm of api] (sched) {scheduler};
  \node[box, below right=6mm and -2mm of api] (cm) {controller\\manager};
  \node[box, below=16mm of api] (kubelet) {kubelet + containerd\\\scriptsize on the node};
  \draw[flow] (api) -- (etcd);
  \draw[flowdash] (sched) -- (api);
  \draw[flowdash] (cm) -- node[lbl,right]{reconcile loops} (api);
  \draw[flow] (kubelet) -- node[lbl,left]{watch \& report} (api);
  \node[lbl, left=8mm of api] (kubectl) {kubectl};
  \draw[flow] (kubectl) -- (api);
\end{tikzpicture}
\caption{Everything talks only to the API server.}
\end{figure}

\begin{description}
  \item[API server] the single front door; \code{kubectl} is a REST
  client for it. Authentication, authorization (Chapter~22's RBAC), and
  validation happen here.
  \item[State store] the desired-state database (etcd in standard
  clusters; k3s embeds sqlite on a single node).
  \item[Scheduler] assigns unscheduled pods to nodes, honoring the
  resource \emph{requests} of Chapter~17.
  \item[Controller manager] the reconciliation loops: Deployment
  $\rightarrow$ ReplicaSet $\rightarrow$ pods, endpoints, and dozens
  more.
  \item[kubelet] the node agent: watches for pods assigned to its node,
  drives containerd to run them, executes the probes, reports status.
\end{description}

k3s packages all of this into one binary --- a CNCF-certified
Kubernetes, minus a VM's overhead, plus batteries: Traefik as Ingress
controller, local-path storage, containerd. Every manifest in this
book would apply unchanged to EKS or OpenShift; that certification is
why learning on k3s transfers.

\section{Objects, labels, namespaces}

Every object is a YAML document with the same skeleton ---
\code{apiVersion}, \code{kind}, \code{metadata} (name, namespace,
labels), \code{spec} (desired), \code{status} (observed, written by
controllers). Two organizing mechanisms:

\begin{description}
  \item[Namespaces] partition the cluster; this project puts everything
  in \code{ts}, keeping it separable from \code{kube-system} and making
  RBAC scopeable --- Jenkins' permissions stop at the namespace edge
  (Chapter~22).
  \item[Labels] arbitrary key/value pairs; \emph{selectors} query them.
  Every wiring in Part~V is a label selector: a Service finds pods by
  label, a Deployment owns pods by label. The kustomization stamps
  \code{app.kubernetes.io/part-of: teacher-supporter} onto everything so
  one selector can list --- or delete --- the whole application.
\end{description}

\section{Getting kubectl to the cluster}

\code{kubectl} reads \code{\textasciitilde/.kube/config}: cluster URL, CA
certificate, and a credential. From Windows, this project reaches
\code{https://100.85.66.43:6443} over Tailscale --- which required
installing k3s with \code{--tls-san 100.85.66.43}, because the API
server's certificate is only valid for the names baked in at install
time.

\begin{handson}[feel the loop]
\begin{lstlisting}[language=bash]
kubectl -n ts get pods
kubectl -n ts delete pod -l app=course-service
kubectl -n ts get pods -w
\end{lstlisting}
Delete the pod, then watch: a replacement appears within seconds,
created not by your command but by the ReplicaSet controller noticing a
missing replica. You never ordered a restart --- you made actual state
diverge from desired, and the loop closed the gap.
\end{handson}

\begin{pitfall}[certificate is valid for\ldots{} not 100.85.66.43]
Symptom: \code{kubectl} from Windows fails TLS validation while working
fine on the server. Cause: the API server certificate lacks the
tailnet address as a SAN --- it was installed without
\code{--tls-san}. Fix at install time is one flag; fix afterwards means
deleting the server TLS directory and restarting k3s. A concrete
instance of a general rule: certificates bind to \emph{names}, so
decide every name a server will be called \emph{before} issuing.
\end{pitfall}

\begin{quiz}
\begin{enumerate}
  \item Explain ``declarative'': what do you write, who acts on it, and
  why does applying the same manifest twice change nothing?
  \item Why is there deliberately no \code{depends\_on} in Kubernetes,
  and what replaces startup ordering during convergence?
  \item Trace \code{kubectl -n ts delete pod ...} through the
  components: who receives it, who notices, who acts, on which machine?
  \item Why did removing the old Ingress from git not remove it from
  the cluster, and what two commands would each have removed it?
\end{enumerate}
\end{quiz}
